\chapter{Experimental Results and Analysis}\label{ch:results}

This chapter presents the empirical results of the retrieval experiments defined in Chapter~\ref{ch:methodology}. The analysis focuses on the relationship between embedding-model capacity and retrieval performance across two proxy tasks. The primary objective is to determine whether increased model capacity yields monotonic performance gains or exhibits saturation and degradation within the constrained context of a small-scale classical music archive.

\section{Experimental Setup Recap}

The experiments were conducted on the full dataset of $N$ tracks from the iPalpiti archive. 
A total of six embedding models were evaluated across two architectural families and three capacity tiers (small, base, and large).

\begin{itemize}
    \item \textbf{CNN-Based Family:} Convolutional architectures spanning compact to high-capacity variants.
    \item \textbf{Transformer-Based Family:} Transformer architectures spanning small to large pretrained checkpoints.
\end{itemize}

Performance is primarily measured using Normalized Discounted Cumulative Gain at rank 10 (NDCG@10) and Precision at rank 10 (P@10). 
Given the relatively small dataset size, interpretation focuses on consistent scaling patterns across capacity tiers rather than marginal numerical differences.


\section{Sanity Proxy Results}

The Sanity Proxy task evaluates the ability of embedding models to retrieve tracks sharing the same composer. This task serves as a baseline for verifying that models capture fundamental compositional structures.

Table~\ref{tab:sanity_results} presents retrieval performance grouped by architectural family and ordered by capacity tier within each family.

\begin{table}[h]
\centering
\caption{Sanity Proxy Results (Composer Retrieval). Metrics reported are corpus-level mean NDCG@10 values.}
\label{tab:sanity_results}
\begin{tabular}{lcc}
\hline
\textbf{Model} & \textbf{Params (M)} & \textbf{NDCG@10} \\
\hline
\multicolumn{3}{c}{\textbf{CNN Family}} \\
CNN-Small & 5 & -- \\
CNN-Base  & 20 & -- \\
CNN-Large & 80 & -- \\
\hline
\multicolumn{3}{c}{\textbf{Transformer Family}} \\
Transformer-Small & 30 & -- \\
Transformer-Base  & 90 & -- \\
Transformer-Large & 330 & -- \\
\hline
\end{tabular}
\end{table}

The observed trend suggests that [TBD: Description of monotonic/non-monotonic trends within families]. Within the CNN family, performance [increases/decreases/plateaus] as capacity increases. Within the Transformer family, performance [increases/decreases/plateaus] relative to parameter growth.

\section{Primary Musical Proxy Results}

The Primary Musical Proxy task evaluates retrieval based on abstract musical character labels (e.g., "Heavy" vs. "Light"). This task represents the core retrieval challenge for the archive.

Table~\ref{tab:musical_results} summarizes the performance across all models.

\begin{table}[h]
\centering
\caption{Primary Musical Proxy Results (Character Retrieval). Metrics reported are corpus-level mean NDCG@10 values.}
\label{tab:musical_results}
\begin{tabular}{lcc}
\hline
\textbf{Model} & \textbf{Params (M)} & \textbf{NDCG@10} \\
\hline
\multicolumn{3}{c}{\textbf{CNN Family}} \\
CNN-Small & 5 & -- \\
CNN-Base  & 20 & -- \\
CNN-Large & 80 & -- \\
\hline
\multicolumn{3}{c}{\textbf{Transformer Family}} \\
Transformer-Small & 30 & -- \\
Transformer-Base  & 90 & -- \\
Transformer-Large & 330 & -- \\
\hline
\end{tabular}
\end{table}

In contrast to the Sanity Proxy, trends in the Musical Proxy task suggest [TBD: Describe difference or similarity in trends]. The relationship between capacity and performance appears to be [stronger/weaker] for this more abstract retrieval task. Notably, the highest-capacity models [do/do not] strictly outperform mid-sized architectures.

\section{Capacity–Performance Trends Within Architectural Families}
This section analyzes scaling behavior strictly within architectural families to isolate the effect of parameter count from architectural inductive bias.
Taken together, the within-family scaling patterns described below provide the primary empirical test of the condition-dependent capacity hypothesis.


\subsection{CNN Scaling Behavior}
Within the CNN family, increasing parameter count from [Model A] to [Model C] results in a [monotonic improvement/saturation/degradation] in NDCG@10. 
Specifically, performance improves from [--] to [--] between the small and base tiers, and from [--] to [--] between the base and large tiers.

\subsection{Transformer Scaling Behavior}
Within the Transformer family, scaling from [Model D] to [Model F] exhibits a [monotonic/non-monotonic] trend. While [Model E] achieves a score of [--], the largest model [Model F] achieves [--]. This observation indicates that for transformer-based architectures in this domain, additional capacity [yields diminishing returns/continues to improve performance/degrades performance], consistent with the hypothesis of condition-dependent scaling.

\section{Cross-Family Descriptive Comparison}

Comparing performance across families reveals distinct behavioral patterns. The best-performing CNN model ([Model Name]) achieves an NDCG@10 of [--], while the best-performing Transformer ([Model Name]) achieves [--].

It is observed that [Compact CNNs/Large Transformers] tend to perform [better/worse/comparably] in the Sanity Proxy task, whereas in the Musical Proxy task, the performance gap [widens/narrows/reverses]. These comparisons are descriptive and do not imply causal superiority of one architecture over another, as parameter count, architecture, and pretraining data scale co-vary between families. However, the data demonstrates that [Low/High] capacity models from the [CNN/Transformer] family are sufficient to achieve competitive retrieval performance in this specific archival setting.

Given the limited scale of the dataset, minor variations in metric values are expected. Therefore, emphasis is placed on consistent trends across capacity tiers rather than isolated numerical differences.

\section{Summary of Empirical Findings}

The experimental results yield the following key observations:

\begin{enumerate}
    \item **Within-Family Scaling:** Increasing model capacity [does/does not] guarantee monotonic performance improvement. Saturation effects were observed in the [CNN/Transformer] family beyond [--] parameters.
    \item **Task Dependence:** Capacity scaling behavior [varies/remains consistent] between the Sanity Proxy and the Musical Proxy tasks.
    \item **Efficiency:** Compact models [competitively match/lag behind] larger models, challenging the assumption that massive parameter counts are strictly necessary for this domain.
\end{enumerate}

Collectively, these findings indicate that embedding-model capacity does not operate as a simple monotonic driver of retrieval quality in small-scale, single-domain archives. Instead, performance patterns vary by task definition and architectural family, supporting a condition-dependent interpretation of scaling behavior.

